\documentclass[uplatex,dvipdfmx,a4paper,11pt]{jsarticle}
\input{../../../templates/kaisetsu.tex}
\kaisetsuprint{統計力学1　第1章「熱力学の復習（熱力学関数・ルジャンドル変換・マクスウェルの関係式）」}{2026-11-11}
%% kaisetsu-moto	problems\統計力学1\shou01\統計力学1-1.tex	kiroku\kaisetsu\統計力学1\統計力学1-1_2026-11-11_komento.tex
\begin{document}
\matome
\begin{saitenichiran}
  \hyougyou{1}{L1・計算}{◎◎◎◎}{S}{全問 ◎}
  \hyougyou{2}{L1・計算}{◎◎◎◎}{S}{全問 ◎}
  \hyougyou{\hyperlink{bt-m-056-maxwell}{3}}{L2・計算}{◎◎◎○}{C}{ヒント2を見た / (4)記述}
  \hyougyou{\hyperlink{bt-m-057-legendre}{4}}{L2・計算}{○○◎◎}{C}{ヒント2を見た / (1)記述 (2)記述}
\end{saitenichiran}
\begin{daiji}
  \item 偏微分は「何を固定したか」までが式の中身である。$(\partial S/\partial V)_{U}=p/T$（基本関係式の微分）と $(\partial S/\partial V)_{T}=(\partial p/\partial T)_{V}$（$F$ から出るマクスウェルの関係式）は別の量で、取り違えると $p=T(\partial p/\partial T)_{V}$ という誤った式になる（問題4 (2)）。
  \item 「確かめよ」の問いは、求めた量を実際に代入して両辺が一致する一行まで書く（問題4 (1) のオイラーの関係式）。
  \item 不可逆過程では、$\Delta U=0$ の理由（熱も仕事も 0）と、$S$ が状態量なので準静的な経路の式の差で $\Delta S$ を求めてよい理由を一言書く（問題3 (4)）。
\end{daiji}
\begin{maemo}
  \item 理由の書き落とし：統計力学1-1（旧版）問題2 (4)（2026-10-27）で「$2^{N}$ になる理由」を書かなかった。今回も問題3 (4) で $\Delta U=0$・$\Delta S$ の求め方の理由が無い。結果を出す式の前に「なぜその式を使ってよいか」を一言置く。
  \item 記号の書き違い：統計力学1-1（旧版）問題4 (3)(4)（2026-10-27）で $S'$ のプライムを落とした。今回は問題4 (2) で $(\partial S/\partial V)_{U}$ を $(\partial p/\partial T)_{V}=(\partial S/\partial V)_{T}$ と書いた。偏微分を書いたら、固定する変数がその式の出どころ（どの関数を何で微分したか）と合っているかを見直す。
\end{maemo}

\mondaihen
\level{2}
\mondaibangou{3}
\begin{mondai}{056-maxwell}{計算}{マクスウェルの関係式による $(\partial U/\partial V)_{T}$ の計算，状態方程式と熱容量からの内部エネルギー・エントロピーの組み立て，自由膨張での温度変化とエントロピーの増加}
物質量 $n$ の気体の状態方程式が，ファン・デル・ワールスの式
\[ p=\frac{nRT}{V-nb}-\frac{an^{2}}{V^{2}} \]
（$a$，$b$ は正の定数）で与えられるとする。体積が十分大きい極限ではこの気体は単原子の理想気体に戻り，定積熱容量は $3nR/2$ になる。
\begin{shoumon}
  \item ギブスの関係式とマクスウェルの関係式を用いて
  \[ \left(\pdv{U}{V}\right)_{T}=T\left(\pdv{p}{T}\right)_{V}-p \]
  を示し，この気体について右辺を計算せよ。
  \item $(\partial C_{V}/\partial V)_{T}$ を $p$ の温度についての2階微分で表し，この気体では $C_{V}=3nR/2$ が任意の $V$ で成り立つことを示せ。さらに，$V\to\infty$ で $U\to3nRT/2$ となるように内部エネルギー $U(T,V)$ を求めよ。
  \item エントロピー $S(T,V)$ を求めよ（加える定数は $S_{0}$ と書いてよい）。
  \item Ar 気体 $\qty{1.00}{\mole}$（$a=\qty{0.1355}{\pascal\metre\tothe{6}\per\mole\squared}$，$b=\qty{3.20e-5}{\cubic\metre\per\mole}$）が，温度 \qty{300}{\kelvin} で体積 $V=\qty{1.00e-3}{\cubic\metre}$ の断熱容器の中にあり，真空の部屋との仕切りを取り除くと全体積 $2V$ に広がった（自由膨張）。膨張後の温度の変化 $\Delta T$ とエントロピーの変化 $\Delta S$ を有効数字3桁で求めよ。$R=\qty{8.314}{\joule\per\kelvin\per\mole}$ とする。
\end{shoumon}

\begin{kotaeran}
  \machigai{4}{○}{記述} $\Delta U=0$ から $\Delta T=\qty{-5.43}{\kelvin}$，$\Delta S=S(T_{2},2V)-S(T_{1},V)=\qty{5.67}{\joule\per\kelvin}$ と，式も数値も正しい。ただし，なぜ $\Delta U=0$ なのか（断熱容器で熱が 0，真空へ広がるので外への仕事が 0）と，自由膨張は不可逆過程なのに (3) の式（準静的な経路で積分したもの）の差で $\Delta S$ を求めてよい理由（$S$ は状態量で，始状態と終状態だけで決まる）が式にも文にも無い。
\end{kotaeran}
\begin{rikai}
  \item $S$ は状態量なので，不可逆過程の $\Delta S$ は，同じ始状態と終状態を結ぶ準静的な経路の $\int\dd S$ で計算できる。自由膨張そのものは準静的ではないので，「$\dd S=\dd Q/T$ で $\dd Q=0$ だから $\Delta S=0$」とはならない。
  \item 自由膨張で保存するのは $U$ であって $T$ ではない（$(\partial U/\partial V)_{T}=an^{2}/V^{2}\neq0$ なので，$U$ 一定のまま体積が増えると温度が下がる）。
\end{rikai}
\end{mondai}

\mondaibangou{4}
\begin{mondai}{057-legendre}{計算}{基本関係式 $S(U,V,N)$ から示強変数・状態方程式を求め，ルジャンドル変換で $F$・$G$ を作ること，オイラーの関係式とギブス--デュエムの関係式}
ある仮想的な物質の基本関係式が
\[ S(U,V,N)=A\,(NVU)^{1/3} \]
（$A$ は正の定数）で与えられるとする。
\begin{shoumon}
  \item $S$ が示量的であることを確かめよ。また，温度 $T$・圧力 $p$・化学ポテンシャル $\mu$ を $U$，$V$，$N$，$S$ で表し，オイラーの関係式 $U=TS-pV+\mu N$ が成り立つことを確かめよ。
  \item 内部エネルギー $U(T,V,N)$，状態方程式 $p(T,V,N)$，ヘルムホルツ自由エネルギー $F(T,V,N)$ を求めよ。
  \item ギブス自由エネルギー $G(T,p,N)$ を求め，$G=\mu N$ を確かめよ。
  \item (3) の $\mu$ を $T$ と $p$ の関数と見て，ギブス--デュエムの関係式から導かれる $(\partial\mu/\partial T)_{p}=-S/N$ と $(\partial\mu/\partial p)_{T}=V/N$ が成り立つことを確かめよ。
\end{shoumon}

\begin{kotaeran}
  \machigai{1}{○}{記述} 示量性の確認と $T=3U/S$，$p=U/V$，$\mu=-U/N$ は正しい。オイラーの関係式は $U=TS-pV+\mu N$ と書いただけで，求めた $T$・$p$・$\mu$ を代入していない。$TS-pV+\mu N=3U-U-U=U$ の一行で確かめが完成する。
  \machigai{2}{○}{記述} $p$ を求める途中で $(\partial p/\partial T)_{V}=(\partial S/\partial V)_{T}=\frac{1}{3}A(NU)^{1/3}V^{-2/3}$ と書いた。右辺は $(\partial S/\partial V)_{U}=p/T$ で，マクスウェルの関係式の $(\partial S/\partial V)_{T}$ ではない。実際 $p=(AT/3)^{3/2}(N/V)^{1/2}$ から $(\partial p/\partial T)_{V}=(3/2)\,p/T$ で，$p/T$ とは違う。次の行の $p=\frac{1}{3}AT\,N^{1/3}V^{-2/3}U^{1/3}$ は $p=T\cdot(p/T)$ として正しく，答えの $p$ と $U$・$F$ も正しい。
\end{kotaeran}
\begin{rikai}
  \item 基本関係式 $S(U,V,N)$ の偏微分は $U,V,N$ のうち残りを固定したもので，$(\partial S/\partial U)_{V,N}=1/T$，$(\partial S/\partial V)_{U,N}=p/T$，$(\partial S/\partial N)_{U,V}=-\mu/T$。状態方程式はこの $p/T$ に $T$ を掛けて，$U$ を $U(T,V,N)$ で消せば得られる。マクスウェルの関係式（$T$ を固定した微分）はここでは要らない。
  \item マクスウェルの関係式 $(\partial S/\partial V)_{T}=(\partial p/\partial T)_{V}$ は，$\dd F=-S\,\dd T-p\,\dd V$ の2階微分の順序の入れ替えから出る。固定する変数（$T$ と $V$）は $F$ の自然な変数と同じ組になっている，と覚えておくと取り違えに気づける。
  \item 確かめ方：$(\partial p/\partial T)_{V}$ と書いた量に $T$ を掛けて $p$ になるのは，$p$ が $T$ の1次式のときだけである。この問題の $p\propto T^{3/2}$ では成り立たない。
\end{rikai}
\end{mondai}

\kaitouhen
\level{2}
\kaisetsubangou{3}
\begin{kaisetsu}{056-maxwell}
\houshin ギブスの関係式から $(\partial U/\partial V)_{T}$ を $(\partial S/\partial V)_{T}$ で書き，マクスウェルの関係式で状態方程式だけで計算できる形にする。$C_{V}$ の体積依存性も状態方程式から決まるので，体積の大きい極限の値と合わせて $U$ と $S$ を組み立てる。自由膨張は $U$ 一定の不可逆過程なので，前後の状態を $U$ の保存で結び，$S$ は状態量として差をとる。
\keisan \shoumonkai{1} $\dd U=T\,\dd S-p\,\dd V$ で $T$ を一定にして $\dd V$ で割ると
\[ \left(\pdv{U}{V}\right)_{T}=T\left(\pdv{S}{V}\right)_{T}-p \]
$\dd F=-S\,\dd T-p\,\dd V$ より $S=-(\partial F/\partial T)_{V}$，$p=-(\partial F/\partial V)_{T}$ で，2階微分の順序を入れ替えると $(\partial S/\partial V)_{T}=(\partial p/\partial T)_{V}$。よって示すべき式が得られる。この気体では $(\partial p/\partial T)_{V}=nR/(V-nb)$ なので
\[ \left(\pdv{U}{V}\right)_{T}=\frac{nRT}{V-nb}-\left(\frac{nRT}{V-nb}-\frac{an^{2}}{V^{2}}\right)=\frac{an^{2}}{V^{2}} \]
\shoumonkai{2} (1) の式を $V$ 一定で $T$ について微分し，偏微分の順序を入れ替える。
\[ \left(\pdv{C_{V}}{V}\right)_{T}=\pdv{}{T}\left[T\left(\pdv{p}{T}\right)_{V}-p\right]_{V}=\left(\pdv{p}{T}\right)_{V}+T\left(\pdv[2]{p}{T}\right)_{V}-\left(\pdv{p}{T}\right)_{V}=T\left(\pdv[2]{p}{T}\right)_{V} \]
$p$ は $T$ の1次式なのでこれは $0$ で，$C_{V}$ は $V$ によらない。$V\to\infty$ での値が $3nR/2$（$T$ によらない）なので，任意の $V$ で $C_{V}=3nR/2$。$U$ は $(\partial U/\partial T)_{V}=3nR/2$ と (1) を満たすので
\[ U(T,V)=\frac{3}{2}nRT-\frac{an^{2}}{V} \]
（$V\to\infty$ で $3nRT/2$ になるので，加える定数は $0$）。
\shoumonkai{3} $\dd S=(\partial S/\partial T)_{V}\dd T+(\partial S/\partial V)_{T}\dd V$ で，$(\partial S/\partial T)_{V}=C_{V}/T$，$(\partial S/\partial V)_{T}=(\partial p/\partial T)_{V}=nR/(V-nb)$。積分して
\[ S(T,V)=\frac{3}{2}nR\log T+nR\log(V-nb)+S_{0} \]
\shoumonkai{4} 自由膨張では外部との熱も仕事も $0$ なので $U$ が保存する。膨張前の温度を $T_{1}=\qty{300}{\kelvin}$，膨張後を $T_{2}$ として
\[ \frac{3}{2}nRT_{1}-\frac{an^{2}}{V}=\frac{3}{2}nRT_{2}-\frac{an^{2}}{2V} \]
より
\[ \Delta T=T_{2}-T_{1}=-\frac{an}{3RV}=-\frac{\num{0.1355}\times\num{1.00}}{3\times\num{8.314}\times\num{1.00e-3}}\,\unit{\kelvin}=\qty{-5.43}{\kelvin} \]
$T_{2}=\qty{294.57}{\kelvin}$ である。エントロピーは状態量なので (3) の式の差をとる。
\[ \Delta S=\frac{3}{2}nR\log\frac{T_{2}}{T_{1}}+nR\log\frac{2V-nb}{V-nb} \]
第1項は $\num{12.47}\times\log(\num{0.98189})\,\unit{\joule\per\kelvin}=\qty{-0.228}{\joule\per\kelvin}$，第2項は
\[ \num{8.314}\times\log\frac{\num{1.968e-3}}{\num{0.968e-3}}\,\unit{\joule\per\kelvin}=\num{8.314}\times\num{0.7095}\,\unit{\joule\per\kelvin}=\qty{5.899}{\joule\per\kelvin} \]
よって $\Delta S=\qty{5.67}{\joule\per\kelvin}$。
\kotae{\shoumonkai{1} $(\partial U/\partial V)_{T}=an^{2}/V^{2}$ \shoumonkai{2} $(\partial C_{V}/\partial V)_{T}=T(\partial^{2}p/\partial T^{2})_{V}=0$，$U=3nRT/2-an^{2}/V$ \shoumonkai{3} $S=\frac{3}{2}nR\log T+nR\log(V-nb)+S_{0}$ \shoumonkai{4} $\Delta T=\qty{-5.43}{\kelvin}$，$\Delta S=\qty{5.67}{\joule\per\kelvin}$}
\begin{ronten}
  \item \shoumonkai{1} $(\partial S/\partial V)_{T}=(\partial p/\partial T)_{V}$ が $F$ の2階微分の順序の入れ替えから出ることを書く。
  \item \shoumonkai{2} $C_{V}$ が $V$ によらないことを示したうえで，体積の大きい極限の値を使って $C_{V}$ と $U$ の定数を決めることを述べる。
  \item \shoumonkai{4} 自由膨張は不可逆過程だが，$S$ は状態量なので，同じ始状態と終状態を結ぶ準静的な経路の式 (3) で差を計算してよいことを述べる。
\end{ronten}
\kenzan 次元：$an^{2}/V^{2}$ は圧力の次元で，$(\partial U/\partial V)_{T}$（エネルギー/体積）と一致する。$an/(RV)$ は $(\unit{\pascal\cubic\metre\per\mole})/(\unit{\joule\per\kelvin\per\mole})=\unit{\kelvin}$。特殊ケース：$a=b=0$ では $\Delta T=0$，$\Delta S=nR\log2=\qty{5.76}{\joule\per\kelvin}$ で理想気体に戻る。$\Delta S>0$ は孤立系の不可逆過程でエントロピーが増えることと整合する。
\tsumazuki $(\partial U/\partial V)_{T}$ を $-p$ としてしまう誤りが多い（それは $(\partial U/\partial V)_{S}$）。自由膨張を準静的断熱過程と取り違えて $\Delta S=0$ としない。
\end{kaisetsu}

\kaisetsubangou{4}
\begin{kaisetsu}{057-legendre}
\houshin 基本関係式の偏微分で示強変数を求め，$1/T$ の式を解いて $U(T,V,N)$ にする。$F=U-TS$，$G=F+pV$ を自然な変数で書き直し，オイラーの関係式・ギブス--デュエムの関係式を具体的に確かめる。
\keisan \shoumonkai{1} $S(\lambda U,\lambda V,\lambda N)=A(\lambda^{3}NVU)^{1/3}=\lambda S(U,V,N)$ なので示量的である。偏微分は
\[ \frac{1}{T}=\pdv{S}{U}=\frac{S}{3U},\qquad \frac{p}{T}=\pdv{S}{V}=\frac{S}{3V},\qquad -\frac{\mu}{T}=\pdv{S}{N}=\frac{S}{3N} \]
よって $T=3U/S$，$p=TS/(3V)=U/V$，$\mu=-TS/(3N)=-U/N$。これらを代入すると
\[ TS-pV+\mu N=3U-U-U=U \]
でオイラーの関係式が成り立つ。
\shoumonkai{2} $1/T=(A/3)(NV)^{1/3}U^{-2/3}$ を $U$ について解くと
\[ U=\left(\frac{AT}{3}\right)^{3/2}(NV)^{1/2} \]
$p=U/V$ より
\[ p=\left(\frac{AT}{3}\right)^{3/2}\left(\frac{N}{V}\right)^{1/2} \]
$TS=3U$ なので
\[ F=U-TS=-2U=-2\left(\frac{AT}{3}\right)^{3/2}(NV)^{1/2} \]
（確認：$-(\partial F/\partial V)_{T,N}=(AT/3)^{3/2}N^{1/2}V^{-1/2}=p$。）
\shoumonkai{3} 状態方程式を $V$ について解くと $V=N(AT/3)^{3}/p^{2}$。$pV=U$ なので $G=F+pV=-2U+U=-U$ で，$U=pV=N(AT/3)^{3}/p$ を使うと
\[ G(T,p,N)=-\frac{N}{p}\left(\frac{AT}{3}\right)^{3} \]
(1) の $\mu=-U/N$ に $U=N(AT/3)^{3}/p$ を入れると $\mu=-(AT/3)^{3}/p$ で，$G=\mu N$ が成り立つ。
\shoumonkai{4} $\mu(T,p)=-(A/3)^{3}T^{3}/p$ を微分する。
\[ \left(\pdv{\mu}{T}\right)_{p}=-\frac{3}{T}\cdot\frac{1}{p}\left(\frac{AT}{3}\right)^{3},\qquad \left(\pdv{\mu}{p}\right)_{T}=\frac{1}{p^{2}}\left(\frac{AT}{3}\right)^{3} \]
一方，$S/N=3U/(NT)=(3/T)(AT/3)^{3}/p$，$V/N=(AT/3)^{3}/p^{2}$ なので，$(\partial\mu/\partial T)_{p}=-S/N$，$(\partial\mu/\partial p)_{T}=V/N$ が成り立つ。これは $N\,\dd\mu=-S\,\dd T+V\,\dd p$（ギブス--デュエムの関係式）を $N$ で割ったものである。
\kotae{\shoumonkai{1} $T=3U/S$，$p=U/V$，$\mu=-U/N$（$TS-pV+\mu N=U$） \shoumonkai{2} $U=(AT/3)^{3/2}(NV)^{1/2}$，$p=(AT/3)^{3/2}(N/V)^{1/2}$，$F=-2(AT/3)^{3/2}(NV)^{1/2}$ \shoumonkai{3} $G=-N(AT/3)^{3}/p$ \shoumonkai{4} （確かめた）}
\begin{ronten}
  \item \shoumonkai{2} $F$ を $T$，$V$，$N$ の関数として書き直してから自然な変数の関数とみなすこと（$U$ や $S$ を残さないこと）を書く。
  \item \shoumonkai{3} $G$ の自然な変数は $T$，$p$，$N$ なので，状態方程式で $V$ を消してから答えることを書く。
\end{ronten}
\kenzan 示量性：$U$，$F$，$G$ はどれも $(V,N)\to(\lambda V,\lambda N)$ で $\lambda$ 倍になり，$p$，$\mu$ は変わらない。$\dd F=-S\,\dd T-p\,\dd V$ と整合（$-(\partial F/\partial T)_{V}=3U/T=S$）。
\tsumazuki $G=\mu N$ を確かめるとき，$\mu$ を $U,V,N$ のまま比べず，同じ変数 $(T,p)$ で書いてから比べる。
\end{kaisetsu}

\end{document}
